\documentclass[../main.tex]{subfiles}
\usepackage[utf8]{inputenc}
\usepackage[T1]{fontenc}
\usepackage[indonesian]{babel}

\begin{document}

\chapter{Kriptografi Modern Simetris: Stream Cipher}

\begin{learningoutcome}
Setelah mempelajari bab ini, mahasiswa diharapkan mampu:
\begin{itemize}
    \item Menjelaskan perbedaan mendasar antara kriptografi klasik dan modern, terutama dalam representasi data dan operasi dasar (Sub-CPMK 1.2).
    \item Memahami konsep operasi bitwise XOR dan aplikasinya dalam algoritma kriptografi (Sub-CPMK 1.2).
    \item Menjelaskan mekanisme kerja \textit{Stream Cipher} dan peran \textit{keystream generator} (Sub-CPMK 1.2).
    \item Menganalisis algoritma \textit{stream cipher} populer seperti RC4, A5/1, Trivium, dan Salsa20/ChaCha20 (Sub-CPMK 1.2).
    \item Membedakan antara \textit{Stream Cipher} dan \textit{Block Cipher} berdasarkan cara pemrosesan data (Sub-CPMK 1.2).
\end{itemize}
\end{learningoutcome}

\section{Pengantar Kriptografi Modern}
Kriptografi modern menandai era setelah penemuan komputer digital, di mana algoritma tidak lagi beroperasi pada level karakter (seperti pada kriptografi klasik), melainkan pada level bit atau byte.

\subsection{Karakteristik Utama}
Perubahan paradigma dari kriptografi klasik ke modern melibatkan beberapa poin kunci:
\begin{itemize}
    \item \textbf{Representasi Biner}: Data (plainteks), kunci, dan cipherteks direpresentasikan sebagai rangkaian bit ($0$ dan $1$).
    \item \textbf{Operasi Dasar}: Operasi \textbf{XOR} ($\oplus$) menjadi operasi yang paling dominan karena efisiensinya dalam perangkat keras dan perangkat lunak.
    \item \textbf{Kompleksitas}: Meskipun tetap menggunakan prinsip substitusi dan transposisi, operasi modern dibuat jauh lebih kompleks melalui putaran (\textit{rounds}), rotasi, dan fungsi non-linear untuk mencegah kriptanalisis.
\end{itemize}

\subsection{Operasi XOR dan Bitwise}
Operasi XOR (\textit{Exclusive OR}) adalah fondasi dari banyak cipher modern.
Tabel kebenaran XOR:
\begin{center}
\begin{tabular}{|c|c|c|}
\hline
$a$ & $b$ & $a \oplus b$ \\
\hline
0 & 0 & 0 \\
0 & 1 & 1 \\
1 & 0 & 1 \\
1 & 1 & 0 \\
\hline
\end{tabular}
\end{center}

Sifat-sifat penting XOR:
\begin{enumerate}
    \item $a \oplus a = 0$
    \item $a \oplus 0 = a$
    \item $a \oplus b = b \oplus a$ (Komutatif)
    \item $a \oplus (b \oplus c) = (a \oplus b) \oplus c$ (Asosiatif)
\end{enumerate}

Sifat $a \oplus a = 0$ memungkinkan proses enkripsi dan dekripsi menggunakan fungsi yang sama:
\[ \text{Enkripsi: } C = P \oplus K \]
\[ \text{Dekripsi: } P = C \oplus K = (P \oplus K) \oplus K = P \oplus (K \oplus K) = P \oplus 0 = P \]

\section{Konsep Stream Cipher (Cipher Alir)}
\textit{Stream cipher} adalah kategori cipher simetris yang mengenkripsi plainteks satu bit atau satu byte pada satu waktu.

\subsection{Mekanisme Kerja}
Inti dari \textit{stream cipher} adalah \textbf{Keystream Generator}. Generator ini menerima kunci rahasia (dan seringkali sebuah \textit{Initialization Vector} atau IV) sebagai \textit{seed} untuk membangkitkan aliran bit acak yang panjangnya sama dengan pesan, yang disebut \textbf{keystream}.

Proses enkripsi dilakukan dengan meng-XOR-kan setiap bit plainteks $p_i$ dengan bit keystream $k_i$ yang berkoresponden:
\[ c_i = p_i \oplus k_i \]

\subsection{Karakteristik Keystream}
Karena tidak mungkin mengirimkan kunci sepanjang pesan, \textit{keystream} dibangkitkan menggunakan algoritma \textit{Pseudo-Random Number Generator} (PRNG). Artinya, aliran bit tersebut terlihat acak, tetapi dapat dibangkitkan ulang secara identik oleh penerima asalkan menggunakan \textit{seed} yang sama.

\subsection{Stream Cipher vs Block Cipher}
\begin{itemize}
    \item \textbf{Stream Cipher}: Memproses data secara kontinu (bit-per-bit/byte-per-byte). Sangat cepat dan cocok untuk data \textit{real-time} (misal: percakapan telepon).
    \item \textbf{Block Cipher}: Memproses data dalam blok berukuran tetap (misal: 128 bit). Lebih kompleks dan umumnya digunakan untuk enkripsi file atau database.
\end{itemize}

\section{Analisis Algoritma Stream Cipher}

\subsection{RC4 (Rivest Cipher 4)}
RC4 adalah salah satu \textit{stream cipher} yang paling populer (digunakan dalam WEP dan SSL/TLS lama). RC4 menggunakan status internal berupa larik $S$ berukuran 256 byte yang dipermutasikan berdasarkan kunci.
\begin{itemize}
    \item \textbf{KSA (Key-Scheduling Algorithm)}: Menginisialisasi dan mengacak larik $S$ berdasarkan kunci.
    \item \textbf{PRGA (Pseudo-Random Generation Algorithm)}: Membangkitkan byte keystream dengan melakukan pertukaran (\textit{swap}) elemen dalam $S$.
\end{itemize}
\textit{Catatan Keamanan}: RC4 kini dianggap tidak aman karena adanya korelasi antara byte-byte awal keystream dengan kunci.

\subsection{A5/1 (GSM Encryption)}
A5/1 digunakan untuk mengamankan komunikasi seluler GSM. Ia menggunakan tiga buah \textit{Linear Feedback Shift Registers} (LFSR) dengan panjang berbeda (19, 22, dan 23 bit). Penentuan apakah sebuah register bergeser atau tidak menggunakan \textbf{kaidah mayoritas} (\textit{majority clocking}).

\subsection{Trivium}
Trivium adalah \textit{lightweight stream cipher} yang dirancang untuk efisiensi tinggi pada perangkat keras. Ia menggunakan tiga buah \textit{Non-Linear Feedback Shift Registers} (NLFSR) yang saling berinteraksi melalui operasi XOR dan AND.

\subsection{Salsa20 dan ChaCha20}
Salsa20 (oleh Daniel J. Bernstein) menggunakan operasi ARX (\textit{Addition, Rotation, XOR}) pada matriks status $4 \times 4$.
\begin{itemize}
    \item \textbf{Salsa20}: Menggunakan \textit{counter} untuk memungkinkan akses acak ke bagian mana pun dari keystream.
    \item \textbf{ChaCha20}: Modifikasi dari Salsa20 untuk meningkatkan difusi. ChaCha20 sangat populer saat ini dan digunakan dalam protokol TLS 1.3.
\end{itemize}

\begin{obeactivity}
\textbf{Aktivitas 5.1: Simulasi XOR Cipher}
1. Tuliskan pesan singkat dalam bentuk biner (8-bit per karakter).
2. Buatlah kunci biner dengan panjang yang sama.
3. Lakukan operasi XOR untuk mendapatkan cipherteks.
4. Berikan cipherteks tersebut kepada rekan Anda dan minta mereka mendekripsinya menggunakan kunci yang sama.
5. Diskusikan apa yang terjadi jika satu bit pada cipherteks diubah (bit-flip).
\end{obeactivity}

\begin{obereflection}
\textbf{Latihan 5.1}
1. Jika plainteks adalah \texttt{10110010} dan kunci adalah \texttt{11001010}, hitunglah cipherteksnya.
2. Mengapa \textit{stream cipher} lebih cocok digunakan untuk komunikasi suara (VoIP) dibandingkan \textit{block cipher}?
3. Jelaskan mengapa penggunaan kunci yang sama untuk dua pesan berbeda dalam \textit{stream cipher} (disebut \textit{key reuse}) sangat berbahaya.
\end{obereflection}

\begin{obeassessment}
\textbf{Kuis Bab 5}
1. Apa perbedaan utama antara \textit{true random} dan \textit{pseudo-random} dalam konteks \textit{keystream generator}?
2. Mengapa operasi XOR dipilih sebagai operasi utama dalam \textit{stream cipher}?
3. Jelaskan prinsip kerja \textit{majority clocking} pada algoritma A5/1.
\end{obeassessment}

\begin{competencychecklist}
\begin{tabular}{|l|c|}
\hline
Indikator Kompetensi & Tercapai \\
\hline
Saya dapat menjelaskan perbedaan kriptografi klasik dan modern & [ ] \\
\hline
Saya dapat melakukan operasi bitwise XOR & [ ] \\
\hline
Saya dapat menjelaskan mekanisme kerja \textit{stream cipher} & [ ] \\
\hline
Saya dapat membedakan karakteristik RC4, A5/1, dan ChaCha20 & [ ] \\
\hline
Saya dapat menganalisis keuntungan \textit{stream cipher} untuk data \textit{real-time} & [ ] \\
\hline
\end{tabular}
\end{center}
\end{competencychecklist}

\end{document}
